% AI-written, AI cross-reviewed; no human sign-off; not a claim that RH is proved
% Atlas unit D570. Local manuscript v3, item G.606, heading on printed PDF p. 308.
% Formalism.LFunctionChainParameters.WeightedSlopeBounds.weighted_mean_lower_bound_le
% paper2/decomposition_tables.tex:5026–5031. Verbatim excerpt; SHA-256 9962b282ee2ca5ab561205aa30c15d39db66554f7eb55bc44e4634989cd4f38c.
% Original TeX uses the macros and cross-reference keys of the source manuscript.

\begin{formalstatement}[D570]\label{decomp:D570}
Under the Hida size setting and Count5, the d(u)-weighted mean of the lower quantities Lo5 at any unit u is at most cSl1·IO(u)+K7v5.
\end{formalstatement}
{\small\raggedright\noindent\textsf{\bfseries Project-specific}\quad \nolinkurl{Formalism.LFunctionChainParameters.WeightedSlopeBounds.weighted_mean_lower_bound_le}\par\noindent \nolinkurl{Formalism/SizeWindowBrackets/ChainSlopeWindows/LevelOfCongruence.lean:162--171}\par\noindent SHA-256 \nolinkurl{5d86d77ec69794c1bbbf19761b4d36c8e6bdd1c717d32688fa94a9fb249443ed}\par\noindent Dominated 12\quad Disjoint 12\quad terminal.\par}
The pointwise lower envelope and the weighted counting constraint are particular to the project’s selected tuple.
\hyperref[lit:S064]{Published sources and search S064}.
